\chapter*{Preface}
\addcontentsline{toc}{chapter}{Preface}

This book exists because of a peculiar experience: spending months trying to
predict something that is, by construction, almost entirely unpredictable ---
and discovering that \emph{almost} is where all the craft lives.

The Jane Street \emph{Real-Time Market Data Forecasting} competition asked
entrants to predict an anonymized financial quantity, \resp{6}, from 79
anonymized features, for roughly forty symbols, at 968 moments of every
trading day, under a streaming protocol that feeds you one timestamp at a
time. The winner's score --- the best score achieved by anyone on Earth, on
years of data, with unlimited compute --- explained \textbf{1.39\%} of the
target's variance. Our own campaign, which this book documents, reached about
1.04\% on a local validation protocol replicating the 8th-place solution, and
produced along the way a body of measurements, derivations, failures, and
tools that generalize to any competition of this shape: DRW's crypto
prediction challenge, Optiver's volatility and order-book contests,
Ubiquant's market forecasting, the M-series forecasting competitions, and the
many private analogues run inside trading firms.

That last clause is the point. \emph{Tabular time-series prediction at very
low signal-to-noise is a genre.} It has its own physics (noise ceilings, drift,
cross-sectional structure), its own failure modes (leakage, unstable
validation, overfit ensembles), its own architecture debates (trees versus
recurrence versus attention), and its own craft knowledge that rarely gets
written down because the people who possess it are either bound by NDAs or
too busy competing. Kaggle write-ups are the genre's samizdat: terse,
brilliant, and cryptic exactly where a beginner needs detail. This book tries
to be the missing long-form treatment --- the document we wish we had found
before starting.

\section*{How this book teaches}

Three commitments shape every chapter:

\begin{enumerate}
\item \textbf{Every principle is paid for with a measurement.} When we claim
  that online refitting matters more than architecture choice, we show the
  ledger: $+0.003$ to $+0.005$ weighted $R^2$ from a daily gradient step,
  versus $\pm 0.0005$ from doubling network width. When we claim optimized
  ensemble weights overfit, we show the optimizer collapsing into corner
  solutions and losing to a plain average. The numbers come from a real
  campaign; the artifact paths in the repository are cited so you can check.
\item \textbf{Negative results get equal billing.} A third of what we know is
  what \emph{didn't} work: training on more history made models worse;
  symbol-identity features were useless; a theoretically beautiful
  path-signature model added nothing an LSTM hadn't already captured.
  Chapter~\ref{ch:negative} collects these; they are as transferable as the
  positive results, and more expensive to rediscover.
\item \textbf{Forensics before modeling.} The single highest-leverage act in
  our campaign was not a model. It was a community member's discovery ---
  reproduced and extended here --- that the target is a moving average of
  near-white noise, shifted forward in time. Everything sensible we did
  afterward flowed from that one piece of reverse engineering. Part~II
  teaches the forensic toolkit as a first-class discipline: autocorrelation
  fingerprints, Fourier signatures, lead--lag maps, deconvolution.
\end{enumerate}

\section*{What you need to bring}

Comfort with probability at the level of ``variance of a sum of independent
variables,'' some exposure to gradient-boosted trees and neural networks, and
working Python. The mathematics is kept honest but grounded: when we derive
the autocorrelation of a moving average or the ridge inverse of a smoothing
operator, it is because the formula is about to earn money in a feature
pipeline, not for its own sake.

\section*{A note on humility}

Everything here was measured on one competition, ported from two others, and
cross-checked where possible. Markets change; protocols differ; the next
contest will break at least one rule in this book. What should transfer is
the \emph{method}: estimate the ceiling, reverse-engineer the process, guard
the information boundary, measure everything twice, and prefer boring
variance reduction over glamorous capacity. Where you find the book wrong,
it will most likely be wrong in a way that the book's own methodology would
have caught. Use the methodology, not the conclusions.

\vspace{2em}
\noindent\textit{Berkeley, July 2026}
